% =====================================================================
%  BUỔI 15 — LỜI GIẢI BÀI TẬP VỀ NHÀ (từ giao-vien/02-loi-giai-bai-tap-ve-nha.md)
%  Sinh tự động bằng scripts/md2tex-luoigiai.py — không sửa tay.
%  Tài liệu giáo viên: dùng để chữa bài trên lớp và soạn đề kiểm tra.
% =====================================================================
\documentclass[11pt, a4paper]{article}
% =====================================================================
%  CAE SHANGHAI — GÓI LỆNH DÙNG CHUNG CHO TÀI LIỆU LUYỆN THI CSCA TOÁN
%  Biên dịch bằng XeLaTeX. Mọi file .tex trong buổi học \input file này.
% =====================================================================

% ---------- NGÔN NGỮ & PHÔNG CHỮ ----------
\usepackage{fontspec}
\usepackage{polyglossia}
\setmainlanguage{vietnamese}
\setotherlanguage{english}

% Chữ Latin (Anh + Việt): TeX Gyre Termes — bản clone Times New Roman,
% đi kèm TeX Live, hỗ trợ đầy đủ dấu tiếng Việt.
\setmainfont{TeX Gyre Termes}
\setsansfont{TeX Gyre Heros}

% Chữ Hán: ưu tiên Source Han Serif SC (思源宋体) — chuẩn học thuật Trung Quốc.
% Nếu máy chưa cài, tự động dùng Songti SC (宋体, có sẵn trên macOS).
% Cả hai đều thuộc họ 宋体 nên diện mạo tài liệu không đổi.
\IfFontExistsTF{Source Han Serif SC}{
  \newcommand{\cjkfamilyname}{Source Han Serif SC}
}{
  \IfFontExistsTF{Songti SC}{
    \newcommand{\cjkfamilyname}{Songti SC}
  }{
    \newcommand{\cjkfamilyname}{Noto Serif CJK SC}
  }
}

% xeCJK tự chuyển font khi gặp ký tự Hán, không cần bọc \cjk{} thủ công.
% Đây là điều kiện bắt buộc: font Latin (TeX Gyre Termes) không có glyph chữ Hán,
% nếu không khai báo xeCJK thì mọi chữ Hán viết trần sẽ bị nuốt khỏi bản in.
\usepackage{xeCJK}
\setCJKfamilyfont{caehan}{\cjkfamilyname}[Scale=0.95]
\setCJKmainfont{\cjkfamilyname}[Scale=0.95]
\setCJKsansfont{\cjkfamilyname}[Scale=0.95]
\setCJKmonofont{\cjkfamilyname}[Scale=0.95]

% \cjk{} ép dùng font Hán trong ngữ cảnh cần thiết (ví dụ trong tiêu đề in đậm
% hoặc trong bảng có định dạng riêng). Dùng \CJKfamily chứ không phải tên font,
% vì tên font chỉ là văn bản — sẽ in ra chữ thay vì đổi font.
\newcommand{\cjk}[1]{{\CJKfamily{caehan}#1}}

% ---------- TOÁN HỌC ----------
\usepackage{amsmath, amssymb, amsthm, mathtools}
\usepackage{unicode-math}
\setmathfont{TeX Gyre Termes Math}

% ---------- BỐ CỤC TRANG ----------
% Lề phải chừa vùng an toàn của letterhead:
%   banner trên chiếm 0–19mm  → top tối thiểu 2,2cm
%   banner dưới chiếm 281–297mm → bottom tối thiểu 2,2cm
% Lề trên 2,6cm (26mm) để tiêu đề đầu trang không đè banner (banner hết ở 19mm).
\usepackage[a4paper, top=2.6cm, bottom=2.4cm, left=2.5cm, right=2.0cm]{geometry}
\usepackage{fancyhdr}
\usepackage{titlesec}
\usepackage{enumitem}
\usepackage{graphicx}
\usepackage{booktabs, array, longtable, tabularx}
\usepackage{xcolor}
\usepackage[most]{tcolorbox}
\usepackage[colorlinks=true, linkcolor=cae_blue, citecolor=cae_blue, urlcolor=cae_blue]{hyperref}

% ---------- BẢNG MÀU THƯƠNG HIỆU ----------
\definecolor{cae_blue}{RGB}{0, 51, 102}
\definecolor{cae_gold}{RGB}{196, 149, 60}
\definecolor{box_bg}{RGB}{245, 247, 250}
\definecolor{box_line}{RGB}{0, 51, 102}

% ---------- TIÊU ĐỀ MỤC ----------
% Đệm đầu trang: banner letterhead chiếm 0–19mm, mà lề trên là 22mm.
% \section* đầu trang (sau \clearpage) sát lề nên phải chừa thêm, nếu không
% sẽ đè lên banner.
\titleformat{\section}
  {\normalfont\large\bfseries\sffamily\color{cae_blue}}{\thesection.}{0.6em}{}
\titleformat{\subsection}
  {\normalfont\normalsize\bfseries\sffamily\color{cae_blue!85!black}}{\thesubsection}{0.6em}{}
\titlespacing*{\section}{0pt}{14pt}{6pt}
\titlespacing*{\subsection}{0pt}{10pt}{4pt}

% Tiêu đề không đánh số ở đầu trang: \section* sau \clearpage nằm sát lề trên
% nên đè lên banner letterhead. Dùng \vspace* ngay TRƯỚC \section* không đủ,
% vì titlesec chèn khoảng cách riêng. Cách chắc chắn là đệm bằng \null.
\newcommand{\sectiondautrang}[1]{\section*{#1}}

% Dạng bài dùng \subsection*{} nên cần định dạng riêng, đậm và có gạch chân màu
\titleformat{\subsubsection}
  {\normalfont\normalsize\bfseries\sffamily\color{cae_blue!85!black}}{\thesubsubsection}{0.6em}{}

% ---------- LETTERHEAD (ảnh nền toàn trang) ----------
% Ảnh letterhead của CAE Shanghai: banner xanh trên/dưới kèm logo chìm giữa.
% Tỉ lệ ảnh 1414x2000 = 0.7070, khớp A4 (0.7071) nên chèn thẳng không méo.
% Đặt ở lớp nền dưới cùng để chữ không bị ảnh che.
% Dùng \AddToShipoutPictureBG (không có dấu *) để áp dụng cho MỌI trang.
\usepackage{eso-pic}
\AddToShipoutPictureBG{%
  \AtPageLowerLeft{%
    \includegraphics[width=\paperwidth,height=\paperheight]{images/letterhead.png}%
  }%
}

% ---------- ĐẦU / CHÂN TRANG ----------
% Letterhead đã có banner trên và dưới, nên bỏ dòng chạy để tránh chồng chữ.
% Số trang đặt trong vùng nội dung, sát trên banner dưới, lề phải.
\pagestyle{fancy}
\setlength{\headheight}{0pt}
\fancyhf{}
\renewcommand{\headrulewidth}{0pt}
\renewcommand{\footrulewidth}{0pt}
\fancyfoot[R]{\raisebox{0.9cm}{\small\color{cae_blue}\thepage}}

% ---------- MÔI TRƯỜNG NỘI DUNG ----------
% Trong tcolorbox, LaTeX đặt parindent = parskip = 0, nên các đoạn văn trong
% khung dính liền nhau. Đặt lại khoảng cách đoạn cho mọi khung.
\newcommand{\khoangcachdoan}{\setlength{\parskip}{0.5em}\setlength{\parindent}{0pt}}
\tcbset{before upper={\khoangcachdoan}}

\newtcolorbox{dinhnghia}[1][]{
  colback=box_bg, colframe=box_line, boxrule=0.6pt, arc=1pt,
  fonttitle=\bfseries, title={Định nghĩa #1}, breakable}
\newtcolorbox{dinhly}[1][]{
  colback=box_bg, colframe=cae_gold, boxrule=0.6pt, arc=1pt,
  fonttitle=\bfseries, title={Tính chất #1}, breakable}
\newtcolorbox{vidu}[1][]{
  colback=white, colframe=black!40, boxrule=0.4pt, arc=1pt,
  fonttitle=\bfseries, title={Ví dụ #1}, breakable}
\newtcolorbox{baitap}[1][]{
  colback=white, colframe=cae_blue!50, boxrule=0.4pt, arc=1pt,
  fonttitle=\bfseries, title={Câu #1}, breakable}

% Ô đáp án: viền xanh đậm hơn, dùng trong file đáp án
\newtcolorbox{dapan}[1][]{
  colback=box_bg, colframe=cae_blue, boxrule=0.6pt, arc=1pt,
  fonttitle=\bfseries, title={Câu #1}, breakable}

% ---------- TÊN CHỨNG MINH ----------
\renewcommand{\proofname}{\normalfont\bfseries Chứng minh}

% ---------- MÔN HỌC ----------
% Đổi được ở file .tex trước khi % =====================================================================
%  CAE SHANGHAI — GÓI LỆNH DÙNG CHUNG CHO TÀI LIỆU LUYỆN THI CSCA TOÁN
%  Biên dịch bằng XeLaTeX. Mọi file .tex trong buổi học \input file này.
% =====================================================================

% ---------- NGÔN NGỮ & PHÔNG CHỮ ----------
\usepackage{fontspec}
\usepackage{polyglossia}
\setmainlanguage{vietnamese}
\setotherlanguage{english}

% Chữ Latin (Anh + Việt): TeX Gyre Termes — bản clone Times New Roman,
% đi kèm TeX Live, hỗ trợ đầy đủ dấu tiếng Việt.
\setmainfont{TeX Gyre Termes}
\setsansfont{TeX Gyre Heros}

% Chữ Hán: ưu tiên Source Han Serif SC (思源宋体) — chuẩn học thuật Trung Quốc.
% Nếu máy chưa cài, tự động dùng Songti SC (宋体, có sẵn trên macOS).
% Cả hai đều thuộc họ 宋体 nên diện mạo tài liệu không đổi.
\IfFontExistsTF{Source Han Serif SC}{
  \newcommand{\cjkfamilyname}{Source Han Serif SC}
}{
  \IfFontExistsTF{Songti SC}{
    \newcommand{\cjkfamilyname}{Songti SC}
  }{
    \newcommand{\cjkfamilyname}{Noto Serif CJK SC}
  }
}

% xeCJK tự chuyển font khi gặp ký tự Hán, không cần bọc \cjk{} thủ công.
% Đây là điều kiện bắt buộc: font Latin (TeX Gyre Termes) không có glyph chữ Hán,
% nếu không khai báo xeCJK thì mọi chữ Hán viết trần sẽ bị nuốt khỏi bản in.
\usepackage{xeCJK}
\setCJKfamilyfont{caehan}{\cjkfamilyname}[Scale=0.95]
\setCJKmainfont{\cjkfamilyname}[Scale=0.95]
\setCJKsansfont{\cjkfamilyname}[Scale=0.95]
\setCJKmonofont{\cjkfamilyname}[Scale=0.95]

% \cjk{} ép dùng font Hán trong ngữ cảnh cần thiết (ví dụ trong tiêu đề in đậm
% hoặc trong bảng có định dạng riêng). Dùng \CJKfamily chứ không phải tên font,
% vì tên font chỉ là văn bản — sẽ in ra chữ thay vì đổi font.
\newcommand{\cjk}[1]{{\CJKfamily{caehan}#1}}

% ---------- TOÁN HỌC ----------
\usepackage{amsmath, amssymb, amsthm, mathtools}
\usepackage{unicode-math}
\setmathfont{TeX Gyre Termes Math}

% ---------- BỐ CỤC TRANG ----------
% Lề phải chừa vùng an toàn của letterhead:
%   banner trên chiếm 0–19mm  → top tối thiểu 2,2cm
%   banner dưới chiếm 281–297mm → bottom tối thiểu 2,2cm
% Lề trên 2,6cm (26mm) để tiêu đề đầu trang không đè banner (banner hết ở 19mm).
\usepackage[a4paper, top=2.6cm, bottom=2.4cm, left=2.5cm, right=2.0cm]{geometry}
\usepackage{fancyhdr}
\usepackage{titlesec}
\usepackage{enumitem}
\usepackage{graphicx}
\usepackage{booktabs, array, longtable, tabularx}
\usepackage{xcolor}
\usepackage[most]{tcolorbox}
\usepackage[colorlinks=true, linkcolor=cae_blue, citecolor=cae_blue, urlcolor=cae_blue]{hyperref}

% ---------- BẢNG MÀU THƯƠNG HIỆU ----------
\definecolor{cae_blue}{RGB}{0, 51, 102}
\definecolor{cae_gold}{RGB}{196, 149, 60}
\definecolor{box_bg}{RGB}{245, 247, 250}
\definecolor{box_line}{RGB}{0, 51, 102}

% ---------- TIÊU ĐỀ MỤC ----------
% Đệm đầu trang: banner letterhead chiếm 0–19mm, mà lề trên là 22mm.
% \section* đầu trang (sau \clearpage) sát lề nên phải chừa thêm, nếu không
% sẽ đè lên banner.
\titleformat{\section}
  {\normalfont\large\bfseries\sffamily\color{cae_blue}}{\thesection.}{0.6em}{}
\titleformat{\subsection}
  {\normalfont\normalsize\bfseries\sffamily\color{cae_blue!85!black}}{\thesubsection}{0.6em}{}
\titlespacing*{\section}{0pt}{14pt}{6pt}
\titlespacing*{\subsection}{0pt}{10pt}{4pt}

% Tiêu đề không đánh số ở đầu trang: \section* sau \clearpage nằm sát lề trên
% nên đè lên banner letterhead. Dùng \vspace* ngay TRƯỚC \section* không đủ,
% vì titlesec chèn khoảng cách riêng. Cách chắc chắn là đệm bằng \null.
\newcommand{\sectiondautrang}[1]{\section*{#1}}

% Dạng bài dùng \subsection*{} nên cần định dạng riêng, đậm và có gạch chân màu
\titleformat{\subsubsection}
  {\normalfont\normalsize\bfseries\sffamily\color{cae_blue!85!black}}{\thesubsubsection}{0.6em}{}

% ---------- LETTERHEAD (ảnh nền toàn trang) ----------
% Ảnh letterhead của CAE Shanghai: banner xanh trên/dưới kèm logo chìm giữa.
% Tỉ lệ ảnh 1414x2000 = 0.7070, khớp A4 (0.7071) nên chèn thẳng không méo.
% Đặt ở lớp nền dưới cùng để chữ không bị ảnh che.
% Dùng \AddToShipoutPictureBG (không có dấu *) để áp dụng cho MỌI trang.
\usepackage{eso-pic}
\AddToShipoutPictureBG{%
  \AtPageLowerLeft{%
    \includegraphics[width=\paperwidth,height=\paperheight]{images/letterhead.png}%
  }%
}

% ---------- ĐẦU / CHÂN TRANG ----------
% Letterhead đã có banner trên và dưới, nên bỏ dòng chạy để tránh chồng chữ.
% Số trang đặt trong vùng nội dung, sát trên banner dưới, lề phải.
\pagestyle{fancy}
\setlength{\headheight}{0pt}
\fancyhf{}
\renewcommand{\headrulewidth}{0pt}
\renewcommand{\footrulewidth}{0pt}
\fancyfoot[R]{\raisebox{0.9cm}{\small\color{cae_blue}\thepage}}

% ---------- MÔI TRƯỜNG NỘI DUNG ----------
% Trong tcolorbox, LaTeX đặt parindent = parskip = 0, nên các đoạn văn trong
% khung dính liền nhau. Đặt lại khoảng cách đoạn cho mọi khung.
\newcommand{\khoangcachdoan}{\setlength{\parskip}{0.5em}\setlength{\parindent}{0pt}}
\tcbset{before upper={\khoangcachdoan}}

\newtcolorbox{dinhnghia}[1][]{
  colback=box_bg, colframe=box_line, boxrule=0.6pt, arc=1pt,
  fonttitle=\bfseries, title={Định nghĩa #1}, breakable}
\newtcolorbox{dinhly}[1][]{
  colback=box_bg, colframe=cae_gold, boxrule=0.6pt, arc=1pt,
  fonttitle=\bfseries, title={Tính chất #1}, breakable}
\newtcolorbox{vidu}[1][]{
  colback=white, colframe=black!40, boxrule=0.4pt, arc=1pt,
  fonttitle=\bfseries, title={Ví dụ #1}, breakable}
\newtcolorbox{baitap}[1][]{
  colback=white, colframe=cae_blue!50, boxrule=0.4pt, arc=1pt,
  fonttitle=\bfseries, title={Câu #1}, breakable}

% Ô đáp án: viền xanh đậm hơn, dùng trong file đáp án
\newtcolorbox{dapan}[1][]{
  colback=box_bg, colframe=cae_blue, boxrule=0.6pt, arc=1pt,
  fonttitle=\bfseries, title={Câu #1}, breakable}

% ---------- TÊN CHỨNG MINH ----------
\renewcommand{\proofname}{\normalfont\bfseries Chứng minh}

% ---------- MÔN HỌC ----------
% Đổi được ở file .tex trước khi % =====================================================================
%  CAE SHANGHAI — GÓI LỆNH DÙNG CHUNG CHO TÀI LIỆU LUYỆN THI CSCA TOÁN
%  Biên dịch bằng XeLaTeX. Mọi file .tex trong buổi học \input file này.
% =====================================================================

% ---------- NGÔN NGỮ & PHÔNG CHỮ ----------
\usepackage{fontspec}
\usepackage{polyglossia}
\setmainlanguage{vietnamese}
\setotherlanguage{english}

% Chữ Latin (Anh + Việt): TeX Gyre Termes — bản clone Times New Roman,
% đi kèm TeX Live, hỗ trợ đầy đủ dấu tiếng Việt.
\setmainfont{TeX Gyre Termes}
\setsansfont{TeX Gyre Heros}

% Chữ Hán: ưu tiên Source Han Serif SC (思源宋体) — chuẩn học thuật Trung Quốc.
% Nếu máy chưa cài, tự động dùng Songti SC (宋体, có sẵn trên macOS).
% Cả hai đều thuộc họ 宋体 nên diện mạo tài liệu không đổi.
\IfFontExistsTF{Source Han Serif SC}{
  \newcommand{\cjkfamilyname}{Source Han Serif SC}
}{
  \IfFontExistsTF{Songti SC}{
    \newcommand{\cjkfamilyname}{Songti SC}
  }{
    \newcommand{\cjkfamilyname}{Noto Serif CJK SC}
  }
}

% xeCJK tự chuyển font khi gặp ký tự Hán, không cần bọc \cjk{} thủ công.
% Đây là điều kiện bắt buộc: font Latin (TeX Gyre Termes) không có glyph chữ Hán,
% nếu không khai báo xeCJK thì mọi chữ Hán viết trần sẽ bị nuốt khỏi bản in.
\usepackage{xeCJK}
\setCJKfamilyfont{caehan}{\cjkfamilyname}[Scale=0.95]
\setCJKmainfont{\cjkfamilyname}[Scale=0.95]
\setCJKsansfont{\cjkfamilyname}[Scale=0.95]
\setCJKmonofont{\cjkfamilyname}[Scale=0.95]

% \cjk{} ép dùng font Hán trong ngữ cảnh cần thiết (ví dụ trong tiêu đề in đậm
% hoặc trong bảng có định dạng riêng). Dùng \CJKfamily chứ không phải tên font,
% vì tên font chỉ là văn bản — sẽ in ra chữ thay vì đổi font.
\newcommand{\cjk}[1]{{\CJKfamily{caehan}#1}}

% ---------- TOÁN HỌC ----------
\usepackage{amsmath, amssymb, amsthm, mathtools}
\usepackage{unicode-math}
\setmathfont{TeX Gyre Termes Math}

% ---------- BỐ CỤC TRANG ----------
% Lề phải chừa vùng an toàn của letterhead:
%   banner trên chiếm 0–19mm  → top tối thiểu 2,2cm
%   banner dưới chiếm 281–297mm → bottom tối thiểu 2,2cm
% Lề trên 2,6cm (26mm) để tiêu đề đầu trang không đè banner (banner hết ở 19mm).
\usepackage[a4paper, top=2.6cm, bottom=2.4cm, left=2.5cm, right=2.0cm]{geometry}
\usepackage{fancyhdr}
\usepackage{titlesec}
\usepackage{enumitem}
\usepackage{graphicx}
\usepackage{booktabs, array, longtable, tabularx}
\usepackage{xcolor}
\usepackage[most]{tcolorbox}
\usepackage[colorlinks=true, linkcolor=cae_blue, citecolor=cae_blue, urlcolor=cae_blue]{hyperref}

% ---------- BẢNG MÀU THƯƠNG HIỆU ----------
\definecolor{cae_blue}{RGB}{0, 51, 102}
\definecolor{cae_gold}{RGB}{196, 149, 60}
\definecolor{box_bg}{RGB}{245, 247, 250}
\definecolor{box_line}{RGB}{0, 51, 102}

% ---------- TIÊU ĐỀ MỤC ----------
% Đệm đầu trang: banner letterhead chiếm 0–19mm, mà lề trên là 22mm.
% \section* đầu trang (sau \clearpage) sát lề nên phải chừa thêm, nếu không
% sẽ đè lên banner.
\titleformat{\section}
  {\normalfont\large\bfseries\sffamily\color{cae_blue}}{\thesection.}{0.6em}{}
\titleformat{\subsection}
  {\normalfont\normalsize\bfseries\sffamily\color{cae_blue!85!black}}{\thesubsection}{0.6em}{}
\titlespacing*{\section}{0pt}{14pt}{6pt}
\titlespacing*{\subsection}{0pt}{10pt}{4pt}

% Tiêu đề không đánh số ở đầu trang: \section* sau \clearpage nằm sát lề trên
% nên đè lên banner letterhead. Dùng \vspace* ngay TRƯỚC \section* không đủ,
% vì titlesec chèn khoảng cách riêng. Cách chắc chắn là đệm bằng \null.
\newcommand{\sectiondautrang}[1]{\section*{#1}}

% Dạng bài dùng \subsection*{} nên cần định dạng riêng, đậm và có gạch chân màu
\titleformat{\subsubsection}
  {\normalfont\normalsize\bfseries\sffamily\color{cae_blue!85!black}}{\thesubsubsection}{0.6em}{}

% ---------- LETTERHEAD (ảnh nền toàn trang) ----------
% Ảnh letterhead của CAE Shanghai: banner xanh trên/dưới kèm logo chìm giữa.
% Tỉ lệ ảnh 1414x2000 = 0.7070, khớp A4 (0.7071) nên chèn thẳng không méo.
% Đặt ở lớp nền dưới cùng để chữ không bị ảnh che.
% Dùng \AddToShipoutPictureBG (không có dấu *) để áp dụng cho MỌI trang.
\usepackage{eso-pic}
\AddToShipoutPictureBG{%
  \AtPageLowerLeft{%
    \includegraphics[width=\paperwidth,height=\paperheight]{images/letterhead.png}%
  }%
}

% ---------- ĐẦU / CHÂN TRANG ----------
% Letterhead đã có banner trên và dưới, nên bỏ dòng chạy để tránh chồng chữ.
% Số trang đặt trong vùng nội dung, sát trên banner dưới, lề phải.
\pagestyle{fancy}
\setlength{\headheight}{0pt}
\fancyhf{}
\renewcommand{\headrulewidth}{0pt}
\renewcommand{\footrulewidth}{0pt}
\fancyfoot[R]{\raisebox{0.9cm}{\small\color{cae_blue}\thepage}}

% ---------- MÔI TRƯỜNG NỘI DUNG ----------
% Trong tcolorbox, LaTeX đặt parindent = parskip = 0, nên các đoạn văn trong
% khung dính liền nhau. Đặt lại khoảng cách đoạn cho mọi khung.
\newcommand{\khoangcachdoan}{\setlength{\parskip}{0.5em}\setlength{\parindent}{0pt}}
\tcbset{before upper={\khoangcachdoan}}

\newtcolorbox{dinhnghia}[1][]{
  colback=box_bg, colframe=box_line, boxrule=0.6pt, arc=1pt,
  fonttitle=\bfseries, title={Định nghĩa #1}, breakable}
\newtcolorbox{dinhly}[1][]{
  colback=box_bg, colframe=cae_gold, boxrule=0.6pt, arc=1pt,
  fonttitle=\bfseries, title={Tính chất #1}, breakable}
\newtcolorbox{vidu}[1][]{
  colback=white, colframe=black!40, boxrule=0.4pt, arc=1pt,
  fonttitle=\bfseries, title={Ví dụ #1}, breakable}
\newtcolorbox{baitap}[1][]{
  colback=white, colframe=cae_blue!50, boxrule=0.4pt, arc=1pt,
  fonttitle=\bfseries, title={Câu #1}, breakable}

% Ô đáp án: viền xanh đậm hơn, dùng trong file đáp án
\newtcolorbox{dapan}[1][]{
  colback=box_bg, colframe=cae_blue, boxrule=0.6pt, arc=1pt,
  fonttitle=\bfseries, title={Câu #1}, breakable}

% ---------- TÊN CHỨNG MINH ----------
\renewcommand{\proofname}{\normalfont\bfseries Chứng minh}

% ---------- MÔN HỌC ----------
% Đổi được ở file .tex trước khi \input{cae-style} bằng \renewcommand{\monhoc}{...}
\providecommand{\monhoc}{Môn Toán}

% ---------- TRANG BÌA GỌN ----------
% #1 = loại tài liệu (TÀI LIỆU CHUẨN BỊ, TÀI LIỆU BUỔI HỌC, ...)
% #2 = số buổi
% #3 = tên buổi tiếng Việt
% #4 = tên buổi tiếng Hán
% #5 = module và tỉ trọng
\newcommand{\trangbia}[5]{%
  \begin{center}
    \vspace*{0.6cm}
    {\large\bfseries\color{cae_blue} CAE SHANGHAI}\\[0.15cm]
    {\footnotesize Tài liệu luyện thi CSCA · \monhoc}\\[0.5cm]
    {\color{cae_gold}\rule{0.55\linewidth}{0.8pt}}\\[0.35cm]
    {\small\bfseries\color{cae_blue!80!black} #1}\\[0.35cm]
    {\Large\bfseries\color{cae_blue} BUỔI #2 — #3}\\[0.2cm]
    {\large\color{cae_blue!85!black} \cjk{#4}}\\[0.25cm]
    \ifx\relax#5\relax\else{\small #5}\\[0.35cm]\fi
    {\color{cae_gold}\rule{0.55\linewidth}{0.8pt}}\\[0.3cm]
    {\footnotesize Biên soạn: Hoàng Dương, CAE Shanghai \quad·\quad \today}
  \end{center}
  \vspace{0.3cm}
}

% ---------- TRANG BÌA CHO ĐỀ KIỂM TRA ----------
% Khác \trangbia ở chỗ không có chữ "BUỔI N" — dùng cho đề thi, không phải bài giảng.
% #1 = loại tài liệu, #2 = tên kỳ thi (Việt), #3 = tên kỳ thi (Hán), #4 = thông số
\newcommand{\biade}[4]{%
  \begin{center}
    \vspace*{0.6cm}
    {\large\bfseries\color{cae_blue} CAE SHANGHAI}\\[0.15cm]
    {\footnotesize Tài liệu luyện thi CSCA · \monhoc}\\[0.5cm]
    {\color{cae_gold}\rule{0.55\linewidth}{0.8pt}}\\[0.35cm]
    {\small\bfseries\color{cae_blue!80!black} #1}\\[0.35cm]
    {\Large\bfseries\color{cae_blue} #2}\\[0.2cm]
    {\large\color{cae_blue!85!black} \cjk{#3}}\\[0.25cm]
    {\small #4}\\[0.35cm]
    {\color{cae_gold}\rule{0.55\linewidth}{0.8pt}}\\[0.3cm]
    {\footnotesize Biên soạn: Hoàng Dương, CAE Shanghai \quad·\quad \today}
  \end{center}
  \vspace{0.3cm}
}

% ---------- TIỆN ÍCH ----------
% Nhãn từ khóa in đậm ở đầu đoạn, dùng cho Từ vựng / Từ khóa nhận dạng / Chú ý
% Khoảng cách đoạn do khung tự đặt (xem \khoangcachdoan ở trên),
% ở đây chỉ cần tách đoạn và bỏ thụt đầu dòng.
\newcommand{\nhan}[1]{\par\noindent\textbf{#1.}\ \ignorespaces}

% Bảng gọn, không đánh số
\newcommand{\bang}[1]{%
  \begin{center}#1\end{center}}

% Đánh dấu ô trống cho học sinh điền
\newcommand{\chotrong}{\makebox[3.2cm]{\dotfill}}
 bằng \renewcommand{\monhoc}{...}
\providecommand{\monhoc}{Môn Toán}

% ---------- TRANG BÌA GỌN ----------
% #1 = loại tài liệu (TÀI LIỆU CHUẨN BỊ, TÀI LIỆU BUỔI HỌC, ...)
% #2 = số buổi
% #3 = tên buổi tiếng Việt
% #4 = tên buổi tiếng Hán
% #5 = module và tỉ trọng
\newcommand{\trangbia}[5]{%
  \begin{center}
    \vspace*{0.6cm}
    {\large\bfseries\color{cae_blue} CAE SHANGHAI}\\[0.15cm]
    {\footnotesize Tài liệu luyện thi CSCA · \monhoc}\\[0.5cm]
    {\color{cae_gold}\rule{0.55\linewidth}{0.8pt}}\\[0.35cm]
    {\small\bfseries\color{cae_blue!80!black} #1}\\[0.35cm]
    {\Large\bfseries\color{cae_blue} BUỔI #2 — #3}\\[0.2cm]
    {\large\color{cae_blue!85!black} \cjk{#4}}\\[0.25cm]
    \ifx\relax#5\relax\else{\small #5}\\[0.35cm]\fi
    {\color{cae_gold}\rule{0.55\linewidth}{0.8pt}}\\[0.3cm]
    {\footnotesize Biên soạn: Hoàng Dương, CAE Shanghai \quad·\quad \today}
  \end{center}
  \vspace{0.3cm}
}

% ---------- TRANG BÌA CHO ĐỀ KIỂM TRA ----------
% Khác \trangbia ở chỗ không có chữ "BUỔI N" — dùng cho đề thi, không phải bài giảng.
% #1 = loại tài liệu, #2 = tên kỳ thi (Việt), #3 = tên kỳ thi (Hán), #4 = thông số
\newcommand{\biade}[4]{%
  \begin{center}
    \vspace*{0.6cm}
    {\large\bfseries\color{cae_blue} CAE SHANGHAI}\\[0.15cm]
    {\footnotesize Tài liệu luyện thi CSCA · \monhoc}\\[0.5cm]
    {\color{cae_gold}\rule{0.55\linewidth}{0.8pt}}\\[0.35cm]
    {\small\bfseries\color{cae_blue!80!black} #1}\\[0.35cm]
    {\Large\bfseries\color{cae_blue} #2}\\[0.2cm]
    {\large\color{cae_blue!85!black} \cjk{#3}}\\[0.25cm]
    {\small #4}\\[0.35cm]
    {\color{cae_gold}\rule{0.55\linewidth}{0.8pt}}\\[0.3cm]
    {\footnotesize Biên soạn: Hoàng Dương, CAE Shanghai \quad·\quad \today}
  \end{center}
  \vspace{0.3cm}
}

% ---------- TIỆN ÍCH ----------
% Nhãn từ khóa in đậm ở đầu đoạn, dùng cho Từ vựng / Từ khóa nhận dạng / Chú ý
% Khoảng cách đoạn do khung tự đặt (xem \khoangcachdoan ở trên),
% ở đây chỉ cần tách đoạn và bỏ thụt đầu dòng.
\newcommand{\nhan}[1]{\par\noindent\textbf{#1.}\ \ignorespaces}

% Bảng gọn, không đánh số
\newcommand{\bang}[1]{%
  \begin{center}#1\end{center}}

% Đánh dấu ô trống cho học sinh điền
\newcommand{\chotrong}{\makebox[3.2cm]{\dotfill}}
 bằng \renewcommand{\monhoc}{...}
\providecommand{\monhoc}{Môn Toán}

% ---------- TRANG BÌA GỌN ----------
% #1 = loại tài liệu (TÀI LIỆU CHUẨN BỊ, TÀI LIỆU BUỔI HỌC, ...)
% #2 = số buổi
% #3 = tên buổi tiếng Việt
% #4 = tên buổi tiếng Hán
% #5 = module và tỉ trọng
\newcommand{\trangbia}[5]{%
  \begin{center}
    \vspace*{0.6cm}
    {\large\bfseries\color{cae_blue} CAE SHANGHAI}\\[0.15cm]
    {\footnotesize Tài liệu luyện thi CSCA · \monhoc}\\[0.5cm]
    {\color{cae_gold}\rule{0.55\linewidth}{0.8pt}}\\[0.35cm]
    {\small\bfseries\color{cae_blue!80!black} #1}\\[0.35cm]
    {\Large\bfseries\color{cae_blue} BUỔI #2 — #3}\\[0.2cm]
    {\large\color{cae_blue!85!black} \cjk{#4}}\\[0.25cm]
    \ifx\relax#5\relax\else{\small #5}\\[0.35cm]\fi
    {\color{cae_gold}\rule{0.55\linewidth}{0.8pt}}\\[0.3cm]
    {\footnotesize Biên soạn: Hoàng Dương, CAE Shanghai \quad·\quad \today}
  \end{center}
  \vspace{0.3cm}
}

% ---------- TRANG BÌA CHO ĐỀ KIỂM TRA ----------
% Khác \trangbia ở chỗ không có chữ "BUỔI N" — dùng cho đề thi, không phải bài giảng.
% #1 = loại tài liệu, #2 = tên kỳ thi (Việt), #3 = tên kỳ thi (Hán), #4 = thông số
\newcommand{\biade}[4]{%
  \begin{center}
    \vspace*{0.6cm}
    {\large\bfseries\color{cae_blue} CAE SHANGHAI}\\[0.15cm]
    {\footnotesize Tài liệu luyện thi CSCA · \monhoc}\\[0.5cm]
    {\color{cae_gold}\rule{0.55\linewidth}{0.8pt}}\\[0.35cm]
    {\small\bfseries\color{cae_blue!80!black} #1}\\[0.35cm]
    {\Large\bfseries\color{cae_blue} #2}\\[0.2cm]
    {\large\color{cae_blue!85!black} \cjk{#3}}\\[0.25cm]
    {\small #4}\\[0.35cm]
    {\color{cae_gold}\rule{0.55\linewidth}{0.8pt}}\\[0.3cm]
    {\footnotesize Biên soạn: Hoàng Dương, CAE Shanghai \quad·\quad \today}
  \end{center}
  \vspace{0.3cm}
}

% ---------- TIỆN ÍCH ----------
% Nhãn từ khóa in đậm ở đầu đoạn, dùng cho Từ vựng / Từ khóa nhận dạng / Chú ý
% Khoảng cách đoạn do khung tự đặt (xem \khoangcachdoan ở trên),
% ở đây chỉ cần tách đoạn và bỏ thụt đầu dòng.
\newcommand{\nhan}[1]{\par\noindent\textbf{#1.}\ \ignorespaces}

% Bảng gọn, không đánh số
\newcommand{\bang}[1]{%
  \begin{center}#1\end{center}}

% Đánh dấu ô trống cho học sinh điền
\newcommand{\chotrong}{\makebox[3.2cm]{\dotfill}}


\begin{document}

\trangbia{LỜI GIẢI BÀI TẬP VỀ NHÀ}{15}{Thống kê}{统计}{Tài liệu giáo viên \quad·\quad 40 câu}

\subsection*{Nhóm A — Thống kê}

\begin{dapan}[1 — Đáp án A]
\textbf{Giải.} Dãy có $11$ giá trị: $3, 3, 2, 3, 6, 3, 10, 3, 6, 3, 2$. Số $3$ xuất hiện $6$ lần, các giá trị khác xuất hiện ít hơn, vậy số mốt là $3$.

Kiểm tra các đáp án còn lại: sắp xếp dãy được $2, 2, 3, 3, 3, 3, 3, 3, 6, 6, 10$ nên số trung vị là $3$, trùng với số mốt, vậy B sai. Tổng bằng $44$ nên số trung bình là $\dfrac{44}{11} = 4$, khác $3$, vậy C và D đều sai.
\end{dapan}

\begin{dapan}[2 — Đáp án A]
\textbf{Giải.} Nhóm 甲 có $5$ giá trị $9, 12, x, 24, 27$ đã xếp tăng dần, nên số trung vị là giá trị thứ ba, tức $x = 15$.

Nhóm 乙 có $5$ giá trị $9, 15, 18, y, 24$ với số trung bình bằng $16{,}8$, vậy tổng bằng

$$16{,}8 \times 5 = 84$$

Suy ra $y = 84 - 9 - 15 - 18 - 24 = 18$.

Vậy $x = 15$ và $y = 18$.
\end{dapan}

\begin{dapan}[3 — Đáp án A]
\textbf{Giải.} Phân phối chuẩn đối xứng qua giá trị trung bình $\mu = 2$, nên hai phần đuôi $X < 2 - a$ và $X > 2 + a$ có xác suất bằng nhau.

Đầu mút $0 = 2 - 2$ và đầu mút $4 = 2 + 2$ đối xứng qua $\mu = 2$. Đề cho $P(X < 4) = 0{,}84$, vậy $P(X > 4) = 1 - 0{,}84 = 0{,}16$.

Do tính đối xứng:

$$P(X \le 0) = P(X > 4) = 0{,}16$$

\textbf{Cách nhanh.} Không cần biết $\sigma$. Chỉ cần nhận ra $0$ và $4$ đối xứng qua $\mu = 2$.
\end{dapan}

\begin{dapan}[4 — Đáp án A]
\textbf{Giải.} Hai bạn có cùng số trung bình nên so sánh trực tiếp độ lệch chuẩn. Độ lệch chuẩn của 甲 là $3{,}5$, nhỏ hơn $4{,}62$ của 乙, nên 甲 có mức dao động thấp hơn và phát huy ổn định hơn.
\end{dapan}

\begin{dapan}[5 — Đáp án B]
\textbf{Giải.} Tính tổng điểm từng bạn:

$$\text{甲} = 85 + 90 + 80 + 75 + 80 = 410$$

$$\text{乙} = 80 + 85 + 90 + 80 + 85 = 420$$

Tổng điểm của 乙 cao hơn.

\textbf{Cách nhanh.} So sánh từng môn để tìm phần chênh lệch: 甲 hơn 乙 ở 语文 ($+5$) và 物理 ($-5$, tức 乙 hơn $5$); 乙 hơn ở 数学 ($+5$) và 英语 ($+10$); 化学 bằng nhau. Chênh lệch ròng là $10$ điểm nghiêng về 乙.
\end{dapan}

\begin{dapan}[6 — Đáp án B]
\textbf{Giải.} Số trung bình:

$$\bar{x} = \frac{7 + 8 + 10 + 12 + 13}{5} = \frac{50}{5} = 10$$

Tổng bình phương các độ lệch:

$$(7-10)^2 + (8-10)^2 + (10-10)^2 + (12-10)^2 + (13-10)^2 = 9 + 4 + 0 + 4 + 9 = 26$$

Phương sai:

$$s^2 = \frac{26}{5} = 5{,}2$$
\end{dapan}

\begin{dapan}[7 — Đáp án B]
\textbf{Giải.} Theo quy tắc thực nghiệm, khoảng $\mu \pm \sigma$ chứa khoảng $68\%$ dữ liệu. Với $\mu = 165$ và $\sigma = 5$, khoảng $160$ đến $170$ chứa khoảng $68\%$ số học sinh, tức \textbf{phần lớn} học sinh.

\textbf{Bẫy.} Đáp án A sai vì còn khoảng $32\%$ học sinh nằm ngoài khoảng này. Đáp án D sai vì số trung bình không phải là giá trị có tần số lớn nhất trong mọi bộ dữ liệu; chỉ với phân phối chuẩn thì ba số trung bình, trung vị, mốt mới trùng nhau, nhưng đề chỉ cho trung bình và độ lệch chuẩn.
\end{dapan}

\begin{dapan}[8 — Đáp án C]
\textbf{Giải.} Đọc dãy nồng độ theo thời gian: $50 \to 60 \to 55 \to 40 \to 30$.

Từ 2021 đến 2022 tăng từ $50$ lên $60$, sau đó giảm liên tục qua các năm còn lại. Vậy xu hướng là tăng trước rồi giảm sau.
\end{dapan}

\begin{dapan}[9 — Đáp án C]
\textbf{Giải.} Tần suất của khoảng $[70, 80)$ là $0{,}4$. Nhân với tổng số học sinh là $100$:

$$0{,}4 \times 100 = 40$$
\end{dapan}

\begin{dapan}[10 — Đáp án C]
\textbf{Giải.} Tổng số học sinh là $2 + 5 + 8 + 7 + 2 = 24$, là số chẵn nên số trung vị là trung bình cộng của giá trị thứ $12$ và thứ $13$. Cộng dồn tần số:

\bang{
\begin{tabularx}{\linewidth}{@{}ccX@{}}
\toprule
\textbf{Điểm} & \textbf{Số học sinh} & \textbf{Vị trí} \\
\midrule
$60$ & 2 & $1$ đến $2$ \\
$70$ & 5 & $3$ đến $7$ \\
$80$ & 8 & $8$ đến $15$ \\
$90$ & 7 & $16$ đến $22$ \\
$100$ & 2 & $23$ đến $24$ \\
\bottomrule
\end{tabularx}}

Cả hai vị trí $12$ và $13$ đều ứng với điểm $80$, vậy số trung vị là $80$.
\end{dapan}

\begin{dapan}[11 — Đáp án C]
\textbf{Giải.} Phép biến đổi là trừ đi $2$ ở mọi giá trị. Theo Tính chất 1, phép cộng (hoặc trừ) cùng một số làm số trung bình thay đổi theo nhưng \textbf{không làm thay đổi phương sai}, vì khoảng cách giữa các giá trị trong dãy được giữ nguyên.

Vậy hai mẫu có cùng phương sai.

\textbf{Chú ý.} Số trung bình, số trung vị và số mốt đều bị dịch đi $2$ đơn vị, nên cả ba đáp án A, B, D đều sai.
\end{dapan}

\begin{dapan}[12 — Đáp án C]
\textbf{Giải.} Trong hai tuần, giá lần lượt là $a$ và $b$.

Bạn 甲 mỗi tuần mua $20$ đồng, tổng tiền là $40$ đồng và tổng lượng mua là $\dfrac{20}{a} + \dfrac{20}{b}$, vậy đơn giá trung bình:

$$m\_1 = \frac{40}{\dfrac{20}{a} + \dfrac{20}{b}} = \frac{2ab}{a + b}$$

Bạn 乙 mỗi tuần mua $6$ cân, tổng lượng là $12$ cân và tổng tiền là $6a + 6b$, vậy đơn giá trung bình:

$$m\_2 = \frac{6a + 6b}{12} = \frac{a + b}{2}$$

Xét hiệu, lưu ý $a, b > 0$ và $a \neq b$ vì giá hai tuần khác nhau:

$$m\_2 - m\_1 = \frac{a+b}{2} - \frac{2ab}{a+b} = \frac{(a+b)^2 - 4ab}{2(a+b)} = \frac{(a-b)^2}{2(a+b)} > 0$$

Vậy $m_2 > m_1$.

\textbf{Nhận xét.} Đây là bất đẳng thức Cauchy quen thuộc: trung bình cộng luôn lớn hơn hoặc bằng trung bình điều hòa. Mua theo lượng cố định cho đơn giá trung bình cao hơn mua theo số tiền cố định.
\end{dapan}

\begin{dapan}[13 — Đáp án D]
\textbf{Giải.} Phân phối chuẩn đối xứng qua $\mu = 2$, nên $1{,}9 = 2 - 0{,}1$ và $2{,}1 = 2 + 0{,}1$ là hai điểm đối xứng qua $\mu$.

Vậy $P(X > 2{,}1) = P(X < 1{,}9) = 0{,}1$, suy ra

$$P(X < 2{,}1) = 1 - 0{,}1 = 0{,}9$$
\end{dapan}

\subsection*{Nhóm B — Xác suất}

\begin{dapan}[14 — Đáp án A]
\textbf{Giải.} Số cách lấy $2$ quả từ $a + b$ quả là $C_{a+b}^2 = \dfrac{(a+b)(a+b-1)}{2}$. Số cách lấy $1$ quả đỏ và $1$ quả trắng là $C_a^1 \cdot C_b^1 = ab$.

Xác suất:

$$P = \frac{ab}{C\_{a+b}^2} = \frac{ab}{\dfrac{(a+b)(a+b-1)}{2}} = \frac{2ab}{(a+b)(a+b-1)}$$

Đối chiếu các đáp án: A đúng. B sai vì mẫu $(a+b)^2$ là số cách lấy có hoàn lại, không phải lấy đồng thời. C sai vì thiếu hệ số $2$ ở tử. D sai vì tử là $a+b$ thay vì $ab$.

\textbf{Chú ý.} Đáp án A và C chỉ khác nhau hệ số $2$, đây là bẫy hay gặp: $C_{a+b}^2$ bằng $\dfrac{(a+b)(a+b-1)}{2}$ chứ không phải $(a+b)(a+b-1)$. Học sinh quên chia $2$ sẽ chọn C.
\end{dapan}

\begin{dapan}[15 — Đáp án A]
\textbf{Giải.} Trong $10$ bóng đèn có $3$ bóng hỏng, vậy có $7$ bóng tốt. Lấy lần lượt không hoàn lại $2$ bóng:

\begin{itemize}[leftmargin=1.4em, itemsep=2pt]
\item Lần thứ nhất lấy được bóng tốt: xác suất $\dfrac{7}{10}$.
\item Lần thứ hai lấy được bóng tốt, biết lần đầu đã lấy một bóng tốt: còn $6$ bóng tốt trong $9$ bóng, xác suất $\dfrac{6}{9}$.
\end{itemize}


Vậy xác suất cả hai đều tốt:

$$\frac{7}{10} \times \frac{6}{9}$$
\end{dapan}

\begin{dapan}[16 — Đáp án B]
\textbf{Giải.} Với hai biến cố độc lập, xác suất cả hai cùng xảy ra bằng tích hai xác suất riêng:

$$P(A \cap B) = P(A) \cdot P(B)$$
\end{dapan}

\begin{dapan}[17 — Đáp án B]
\textbf{Giải.} Dùng công thức xác suất toàn phần theo hai loại sản phẩm:

$$P(\text{đạt chuẩn}) = 0{,}4 \times 0{,}95 + 0{,}6 \times 0{,}90 = 0{,}38 + 0{,}54 = 0{,}92$$
\end{dapan}

\begin{dapan}[18 — Đáp án C]
\textbf{Giải.} Không gian mẫu là $\{1, 2, 3, 4, 5, 6\}$.

Biến cố $A = \{1, 4\}$ và $B = \{1, 3, 5\}$ có phần tử chung là $1$, nên $A$ và $B$ \textbf{không} xung khắc và cũng không đối nhau; vậy A và B đều sai.

Hợp của hai biến cố là $A \cup B = \{1, 3, 4, 5\}$, có $4$ phần tử:

$$P(A + B) = \frac{4}{6} = \frac{2}{3}$$
\end{dapan}

\begin{dapan}[19 — Đáp án D]
\textbf{Giải.} Sắp xếp $5$ chữ 仁义礼智信 thành một hàng có $5! = 120$ cách, đều có khả năng như nhau.

Số cách để 仁, 义, 礼 giữ nguyên thứ tự tương đối: trong một hoán vị bất kỳ của $5$ chữ, ba chữ này có $3! = 6$ thứ tự có thể, và chỉ một thứ tự là 仁义礼, nên số cách thỏa mãn là

$$\frac{5!}{3!} = \frac{120}{6} = 20$$

Xác suất:

$$P = \frac{20}{120} = \frac{1}{6}$$

\textbf{Cách nhanh.} Do tính đối xứng, ba chữ 仁, 义, 礼 có $3! = 6$ thứ tự tương đối đồng khả năng, chỉ một thứ tự đúng, nên xác suất là $\dfrac{1}{6}$.
\end{dapan}

\subsection*{Nhóm C — Vector và số phức}

\begin{dapan}[20 — Đáp án B]
\textbf{Giải.} Nhân tử và mẫu với số phức liên hợp của mẫu là $1 + i$:

$$z = \frac{2+i}{1-i} = \frac{(2+i)(1+i)}{(1-i)(1+i)} = \frac{2 + 2i + i + i^2}{1 - i^2}$$

Dùng $i^2 = -1$:

$$z = \frac{2 + 3i - 1}{1 + 1} = \frac{1 + 3i}{2} = \frac{1}{2} + \frac{3}{2}i$$

Số phức liên hợp:

$$\overline{z} = \frac{1}{2} - \frac{3}{2}i$$
\end{dapan}

\begin{dapan}[21 — Đáp án C]
\textbf{Giải.} Khai triển từng phần:

$$(3 - i) \cdot i = 3i - i^2 = 3i + 1$$

Trừ tiếp $(1 + 2i)$:

$$(1 + 3i) - (1 + 2i) = i$$

Vậy kết quả là $i$.
\end{dapan}

\begin{dapan}[22 — Đáp án D]
\textbf{Giải.} Với $a = (1, -5)$ và $b = (2, 4)$:

$$a + b = (3, -1), \qquad a - b = (-1, -9)$$

Tích vô hướng:

$$(a+b) \cdot (a-b) = 3 \cdot (-1) + (-1) \cdot (-9) = -3 + 9 = 6$$

\textbf{Cách nhanh.} Dùng hằng đẳng thức $(a+b) \cdot (a-b) = \lvert a \rvert^2 - \lvert b \rvert^2$:

$$\lvert a \rvert^2 = 1 + 25 = 26, \qquad \lvert b \rvert^2 = 4 + 16 = 20$$

Hiệu bằng $26 - 20 = 6$. Cách này tránh được phép tính tọa độ của hai vector tổng và hiệu.
\end{dapan}

\begin{dapan}[23 — Đáp án D]
\textbf{Giải.} Gọi $\alpha$ là góc giữa hai đường thẳng. Công thức tính cosin góc giữa hai vector chỉ phương:

$$\cos\alpha = \frac{\lvert d\_1 \cdot d\_2 \rvert}{\lvert d\_1 \rvert \cdot \lvert d\_2 \rvert}$$

Tích vô hướng:

$$d\_1 \cdot d\_2 = 1 \cdot 2 + (-1) \cdot 2 + 2 \cdot 1 = 2 - 2 + 2 = 2$$

Độ dài hai vector:

$$\lvert d\_1 \rvert = \sqrt{1 + 1 + 4} = \sqrt{6}, \qquad \lvert d\_2 \rvert = \sqrt{4 + 4 + 1} = 3$$

Vậy

$$\cos\alpha = \frac{2}{3\sqrt{6}} = \frac{2\sqrt{6}}{18} = \frac{\sqrt{6}}{9}$$

\textbf{Chú ý.} Công thức góc giữa hai \textbf{đường thẳng} dùng giá trị tuyệt đối ở tử vì góc giữa hai đường thẳng không vượt quá $90^\circ$. Nếu đề hỏi góc giữa hai \textbf{vector} thì giữ nguyên dấu.
\end{dapan}

\subsection*{Nhóm D — Ôn tập tích lũy}

\begin{dapan}[24 — Đáp án D]
\textbf{Giải.} Xét $x < 0$: biểu thức $x + \dfrac{1}{x}$ nhận giá trị âm và không bị chặn dưới. Khi $x \to 0^-$ thì $\dfrac{1}{x} \to -\infty$, nên biểu thức tiến tới $-\infty$.

Vậy trên tập xác định $\mathbb{R} \setminus \{0\}$, biểu thức \textbf{không có} giá trị nhỏ nhất.

\textbf{Bẫy.} Đáp án B ứng với trường hợp đề chỉ xét $x > 0$, khi đó $x + \dfrac{1}{x} \ge 2$ theo bất đẳng thức Cauchy. Đề không nêu điều kiện nên phải xét cả $x < 0$.
\end{dapan}

\begin{dapan}[25 — Đáp án D]
\textbf{Giải.} Chia cả tử và mẫu cho $\cos\alpha$ (lưu ý $\cos\alpha \neq 0$ vì $\tan\alpha = 2$):

$$\frac{\sin\alpha - \cos\alpha}{\sin\alpha + \cos\alpha} = \frac{\tan\alpha - 1}{\tan\alpha + 1} = \frac{2 - 1}{2 + 1} = \frac{1}{3}$$

\textbf{Cách nhanh.} Gặp biểu thức đẳng cấp bậc nhất theo $\sin$ và $\cos$, luôn chia cả tử và mẫu cho $\cos$ để đưa về $\tan$. Không cần tính $\sin\alpha$ và $\cos\alpha$ riêng.
\end{dapan}

\subsection*{Nhóm D — Ôn tập tích lũy}

\begin{dapan}[1 — Đáp án A]
\textbf{Giải.} Xếp $20$ tuổi theo thứ tự không giảm: $28$; bốn giá trị $29$; hai giá trị $30$; năm giá trị $32$; ba giá trị $36$; ba giá trị $40$; hai giá trị $45$.

Cộng dồn: giá trị thứ $1$ là $28$; các vị trí $2$ đến $5$ là $29$; vị trí $6$ và $7$ là $30$; vị trí $8$ đến $12$ là $32$; vị trí $13$ đến $15$ là $36$; vị trí $16$ đến $18$ là $40$; vị trí $19$ và $20$ là $45$.

Với $20$ phần tử, bách phân vị $25\%$ là trung bình cộng của giá trị thứ $5$ và thứ $6$, tức $\dfrac{29 + 30}{2} = 29{,}5$. Vậy A đúng.

Kiểm tra các đáp án khác: bách phân vị $75\%$ là trung bình cộng của giá trị thứ $15$ và thứ $16$, tức $\dfrac{36 + 40}{2} = 38$, nên B sai. Số trung vị là trung bình cộng của giá trị thứ $10$ và thứ $11$, đều bằng $32$, nên C sai. Số mốt là $32$ (xuất hiện $5$ lần) chứ không phải $5$, nên D sai.
\end{dapan}

\begin{dapan}[2 — Đáp án A]
\textbf{Giải.} Tổng số học sinh toàn trường là $400 + 300 + 200 = 900$. Tỉ lệ học sinh khối 12 (高三) là $\dfrac{200}{900}$.

Số học sinh khối 12 được chọn theo phương pháp chọn mẫu phân tầng:

$$90 \times \frac{200}{900} = 20$$
\end{dapan}

\begin{dapan}[3 — Đáp án A]
\textbf{Giải.} Đường cong phân phối chuẩn đối xứng qua đường thẳng $x = \mu$, cao nhất tại tâm và thấp dần về hai phía, tạo hình chuông.
\end{dapan}

\begin{dapan}[4 — Đáp án A]
\textbf{Giải.} Tính tỉ lệ ủng hộ của từng nhóm:

$$\frac{120}{200} = 0{,}600, \quad \frac{150}{300} = 0{,}500, \quad \frac{100}{200} = 0{,}500, \quad \frac{80}{150} \approx 0{,}533$$

Tỉ lệ cao nhất là $0{,}600$, thuộc nhóm $[18, 30)$.

\textbf{Bẫy.} Nhóm $[30,45)$ có số người ủng hộ đông nhất ($150$) nhưng tỉ lệ chỉ $0{,}5$. Đề hỏi 支持率 (tỉ lệ ủng hộ) chứ không hỏi số người.
\end{dapan}

\begin{dapan}[5 — Đáp án B]
\textbf{Giải.} Dãy có $11$ giá trị với số trung bình bằng $6$, vậy tổng bằng

$$6 \times 11 = 66$$

Tổng mười giá trị đã biết:

$$2 + 4 + 4 + 5 + 5 + 6 + 7 + 8 + 9 + 11 = 61$$

Suy ra $x = 66 - 61 = 5$.

Tổng bình phương các độ lệch so với $6$:

$$(2-6)^2 + (4-6)^2 + (4-6)^2 + (5-6)^2 + (5-6)^2 + (6-6)^2 + (7-6)^2 + (8-6)^2 + (9-6)^2 + (11-6)^2 + (5-6)^2$$

$$= 16 + 4 + 4 + 1 + 1 + 0 + 1 + 4 + 9 + 25 + 1 = 66$$

Phương sai:

$$s^2 = \frac{66}{11} = 6$$
\end{dapan}

\begin{dapan}[6 — Đáp án B]
\textbf{Giải.} Phân phối chuẩn đối xứng qua $\mu = 2$, nên hai điểm $0 = 2 - 2$ và $4 = 2 + 2$ đối xứng qua $\mu$.

Vậy $P(X \le 0) = P(X \geq 4) = 0{,}45$, suy ra

$$P(X \geq 0) = 1 - 0{,}45 = 0{,}55$$
\end{dapan}

\begin{dapan}[7 — Đáp án B]
\textbf{Giải.} Mức ý nghĩa $\alpha$ là xác suất mắc sai lầm loại một, tức xác suất bác bỏ giả thuyết đúng. Với $\alpha = 0{,}05$, xác suất này là $5\%$.
\end{dapan}

\begin{dapan}[8 — Đáp án B]
\textbf{Giải.} Phân phối chuẩn $N(1, \sigma^2)$ đối xứng qua $\mu = 1$. Điều kiện $P(\xi \le 0) = P(\xi \ge a)$ nghĩa là $0$ và $a$ đối xứng qua $1$, suy ra $a = 2$.

Cần tìm giá trị nhỏ nhất của $\dfrac{1}{x} + \dfrac{9}{a - x}$ với $0 < x < a = 2$. Áp dụng hệ quả bất đẳng thức Cauchy cho hai số dương:

$$\frac{1}{x} + \frac{9}{2-x} \ge \frac{(1 + 3)^2}{x + (2-x)} = \frac{16}{2} = 8$$

Dấu bằng xảy ra khi $\dfrac{1}{x} = \dfrac{3}{2-x}$, tức $x = \dfrac{1}{2}$. Vậy giá trị nhỏ nhất là $8$.

\textbf{Cách nhanh.} Với biểu thức dạng $\dfrac{p}{x} + \dfrac{q}{a-x}$ trên $(0, a)$, giá trị nhỏ nhất bằng $\dfrac{(\sqrt{p} + \sqrt{q})^2}{a}$.
\end{dapan}

\begin{dapan}[9 — Đáp án C]
\textbf{Giải.} Từ $N(102, 4^2)$ suy ra $\mu = 102$ và $\sigma = 4$. Đầu mút $114 = 102 + 12 = \mu + 3\sigma$.

Khoảng $\mu \pm 3\sigma$ chứa $99{,}7\%$ dữ liệu, nên phần nằm ngoài khoảng này ở cả hai phía là

$$100\% - 99{,}7\% = 0{,}3\%$$

Do đường cong đối xứng, phần đuôi bên phải bằng nửa phần ngoài khoảng:

$$\frac{0{,}3\%}{2} = 0{,}15\%$$

Kết quả làm tròn theo ba chữ số có nghĩa là $0{,}135\%$.
\end{dapan}

\begin{dapan}[10 — Đáp án C]
\textbf{Giải.} So sánh hàm mật độ đã cho với dạng tổng quát của phân phối chuẩn:

$$f(x) = \frac{1}{\sigma\sqrt{2\pi}} e^{-\dfrac{(x-\mu)^2}{2\sigma^2}}$$

Đề cho $f(x) = \dfrac{1}{\sqrt{8\pi}} e^{-\dfrac{x^2}{8}}$. Từ mẫu số của phân số phía trước suy ra $\sigma\sqrt{2\pi} = \sqrt{8\pi}$, vậy $\sigma = 2$.

Từ số hạng mũ, $(x - \mu)^2 = x^2$ nên $\mu = 0$.

Vậy giá trị trung bình là $0$ và độ lệch chuẩn là $2$.

\textbf{Chú ý.} Trong biểu thức $e^{-\frac{x^2}{8}}$, mẫu số $8$ bằng $2\sigma^2$, không phải $\sigma^2$. Nhầm chỗ này sẽ ra $\sigma = \sqrt{8}$, một đáp án nhiễu.
\end{dapan}

\begin{dapan}[11 — Đáp án C]
\textbf{Giải.} Tỉ lệ đạt chuẩn là $98\%$. Số sản phẩm đạt chuẩn ước lượng trong $100$ sản phẩm:

$$100 \times 0{,}98 = 98$$
\end{dapan}

\begin{dapan}[12 — Đáp án C]
\textbf{Giải.} Hai biến cố "trúng đặc biệt" và "trúng nhất đẳng" xung khắc (một vé chỉ trúng một giải). Tổng số vé trúng một trong hai giải là $1 + 2 = 3$ trên $100$ vé:

$$P = \frac{3}{100} = 0{,}03$$
\end{dapan}

\begin{dapan}[13 — Đáp án D]
\textbf{Giải.} Tổng số học sinh là $150 + 120 + 50 + 30 = 350$. Số nam sinh có thị lực bình thường là $150$.

Xác suất chọn được một nam sinh thị lực bình thường:

$$\frac{150}{350}$$
\end{dapan}

\begin{dapan}[14 — Đáp án D]
\textbf{Giải.} Trong điều tra thống kê, mẫu được gọi là có tính đại diện (代表性) khi nó phản ánh trung thực cấu trúc của tổng thể. Đây là tiêu chuẩn quan trọng nhất khi chọn mẫu.
\end{dapan}

\begin{dapan}[15 — Đáp án D]
\textbf{Giải.} Với phân phối chuẩn $N(\mu, \sigma^2)$, trung bình, trung vị và mốt đều bằng $\mu$. Xét bốn mệnh đề với $a$ là một giá trị bất kỳ:

\begin{itemize}[leftmargin=1.4em, itemsep=2pt]
\item \textbf{甲:} $P(\xi < a - 1) > P(\xi > a + 2)$. Vì $a - 1$ gần tâm hơn $a + 2$ nên xác suất đuôi trái nhỏ hơn, mệnh đề sai.
\item \textbf{乙 và 丙:} $P(\xi > a) = 0{,}5$ và $P(\xi \le a) = 0{,}5$ chỉ đúng khi $a = \mu$, không đúng với mọi $a$. Tuy nhiên đề chỉ cho một mệnh đề sai, nên phải đọc theo nghĩa "tồn tại $a = \mu$ thỏa mãn"; cả hai mệnh đề này được coi là đúng.
\item \textbf{丁:} $P(a < \xi < a + 1) < P(a + 1 < \xi < a + 2)$. Hai khoảng có cùng độ dài nhưng khoảng thứ hai xa tâm hơn, nên xác suất nhỏ hơn, tức phải có dấu $>$. Mệnh đề sai.
\end{itemize}


Vậy mệnh đề sai là 丁.

\textbf{Chú ý.} Quy tắc cần nhớ: với hai đoạn cùng độ dài, đoạn nào \textbf{gần giá trị trung bình hơn} thì chứa xác suất \textbf{lớn hơn}.
\end{dapan}

\clearpage
\sectiondautrang{Tổng kết lỗi thường gặp}
\addcontentsline{toc}{section}{Tổng kết lỗi thường gặp}

Bảng chẩn đoán dùng khi chữa bài và khi soạn đề kiểm tra. Cột ``Câu liên quan'' trỏ tới đề bài tập về nhà.

\bang{
\begin{tabular}{@{}p{5.6cm}p{4.2cm}p{5.2cm}@{}}
\toprule
\textbf{Lỗi} & \textbf{Câu liên quan} & \textbf{Cách cảnh báo học sinh} \\
\midrule
Lấy trung bình cộng của các giá trị, bỏ qua cột tần số & Câu 6, 10 phần chuẩn bị; Câu 6 phần tự chọn & Bảng có cột số người thì phải nhân giá trị với tần số rồi mới chia \\
Tìm trung vị trên dãy chưa sắp xếp & Câu 7 phần chuẩn bị; Câu 2, 10 phần bắt buộc & Viết chữ "sắp xếp" lên đầu bài làm trước khi tính \\
Nhầm 中位数 (trung vị) với 众数 (mốt) & Câu 5 phần bắt buộc; Câu 1 phần tự chọn & Trung vị là giá trị đứng giữa, mốt là giá trị xuất hiện nhiều nhất \\
Nghĩ rằng cộng thêm hằng số làm phương sai thay đổi & Câu 3 phần bắt buộc; Câu 11 phần tự chọn & Cộng thêm chỉ dịch chuyển dãy, khoảng cách giữa các giá trị không đổi \\
Quên chia đôi phần đuôi khi tính tỉ lệ ngoài khoảng $\mu \pm k\sigma$ & Câu 9 phần chuẩn bị; Câu 9 phần tự chọn & Ngoài khoảng thì chia hai, mỗi đuôi chỉ lấy một nửa \\
Nhầm mẫu số $2\sigma^2$ trong hàm mật độ chuẩn & Câu 10 phần tự chọn & Đọc kỹ: mẫu số của số hạng mũ là $2\sigma^2$, không phải $\sigma^2$ \\
Chia tỉ lệ theo số người thay vì theo tổng số của từng nhóm & Câu 4 phần tự chọn & Đề hỏi 支持率 (tỉ lệ) thì lấy từng nhóm chia cho tổng số của chính nhóm đó \\
Không xét đủ trường hợp khi tìm giá trị lớn nhất, nhỏ nhất & Câu 24 phần bắt buộc & Đề không nêu điều kiện của biến thì phải xét cả giá trị âm \\
Quên dấu giá trị tuyệt đối khi tính góc giữa hai đường thẳng & Câu 23 phần bắt buộc & Góc giữa hai đường thẳng không vượt quá $90^\circ$ nên tử phải lấy trị tuyệt đối \\
\bottomrule
\end{tabular}}

% Không dùng \vfill: nó đẩy chữ xuống sát đáy trang,
% nơi banner letterhead chiếm chỗ (281-297mm) và sẽ che mất chữ.
\vspace{0.8cm}
\begin{center}
{\footnotesize\color{cae_blue} Tài liệu nội bộ, CAE SHANGHAI}
\end{center}

\end{document}
